\documentclass[aspectratio=169,10pt]{beamer}
\usetheme{Boadilla}
\usecolortheme{default}
\usepackage[T1]{fontenc}
\usepackage{lmodern}
\usepackage{amsmath,amssymb,booktabs,tabularx,array}
\definecolor{navy}{RGB}{22,48,75}
\definecolor{teal}{RGB}{0,111,112}
\definecolor{muted}{RGB}{90,102,111}
\definecolor{statusred}{RGB}{166,37,43}
\definecolor{statusyellow}{RGB}{255,231,105}
\definecolor{statusamber}{RGB}{238,174,54}
\definecolor{statusgreen}{RGB}{31,122,75}
\setbeamercolor{structure}{fg=teal}
\setbeamercolor{frametitle}{fg=navy}
\setbeamercolor{title}{fg=navy}
\setbeamercolor{footline}{fg=muted}
\setbeamertemplate{navigation symbols}{}
\setbeamertemplate{headline}{\leavevmode\hbox to \paperwidth{\hfill{\tiny\color{muted}ISE 631 \quad Adelman research}\hspace{1.5em}}\vskip2pt}
\setbeamertemplate{footline}{\leavevmode\hbox to \paperwidth{\hspace{1.5em}{\tiny\color{muted}Adelman Project IPR 2}\hfill{\tiny\color{muted}Herbert Jockheck}\hfill{\tiny\color{muted}\insertframenumber/\inserttotalframenumber}\hspace{1.5em}}\vskip4pt}
\setbeamertemplate{itemize items}[circle]
\setlength{\tabcolsep}{4pt}
\newcommand{\src}[1]{\par\vfill{\tiny\color{muted}Source: #1\par}}
\newcommand{\notation}[1]{\par\vspace{0.4em}{\tiny\color{muted}\textbf{Variables:} #1\par}}
\newcommand{\decision}[1]{\par\vspace{0.5em}{\color{teal}\(\bigstar\)\ \textbf{Decision:} #1\par}}
\newcommand{\blocked}{\colorbox{statusred}{\strut\color{white}Blocked}}
\newcommand{\draft}{\colorbox{statusyellow}{\strut\color{black}Draft}}
\newcommand{\notstarted}{\colorbox{statusamber}{\strut\color{black}Not started}}
\newcommand{\complete}{\colorbox{statusgreen}{\strut\color{white}Complete}}
\newcommand{\status}[1]{{\color{teal}\textbf{#1}}}
\newcommand{\R}{\mathcal R}
\newcommand{\U}{\mathcal U}
\newcommand{\E}{\mathbb E}
\title{Dynamic Bid Prices in Revenue Management}
\subtitle{ISE 631 IPR 2 decision brief: theory and evidence plan}
\author{Herbert Jockheck}
\date{Fall 2026}

\begin{document}

% Main slide 1
\begin{frame}[t]{IPR 2 Decision Brief}
\small
\begin{beamercolorbox}[sep=0.7em,center,wd=\textwidth]{block body}
\textbf{Purpose:} Confirm the Adelman mathematics, course-requirement coverage, and evidence plan for the final presentation.
\end{beamercolorbox}
\vspace{0.8em}
\begin{columns}[T,onlytextwidth]
\column{0.31\textwidth}
\textbf{Decisions requested}\par
\begin{itemize}\setlength{\itemsep}{0.5em}
\item \textbf{Slide 8:} Proof and LP depth?
\item \textbf{Slide 10:} IPR 3 computation scope?
\item \textbf{Slide 11:} Published replication target?
\item \textbf{Slide 13:} Conceptual application scope?
\end{itemize}
\column{0.32\textwidth}
\textbf{Agenda today}\par
\begin{itemize}\setlength{\itemsep}{0.5em}
\item Progress and requirement crosswalk
\item Exact, fluid, and affine models
\item Bound proof and column generation
\item Verification and experiment design
\item Final-brief scope and risks
\end{itemize}
\column{0.32\textwidth}
\textbf{Successful end state}\par
\begin{itemize}\setlength{\itemsep}{0.5em}
\item Theory coverage agreed
\item Numerical gates defined
\item IPR 3 comparisons approved
\item Final narrative ordered
\item Extension limits recorded
\end{itemize}
\end{columns}
\src{Structure follows the IPR 1 decision brief, p. 1. Status and requests are project-specific.}
\end{frame}

% Main slide 2
\begin{frame}[t]{Final-presentation crosswalk and current status}
\scriptsize
\begin{tabularx}{\textwidth}{@{}p{0.21\textwidth}p{0.31\textwidth}p{0.16\textwidth}X@{}}
\toprule
Section & Deliverable & Status & Next action\\
\midrule
1. Intro (I14) & Pricing motivation and approved RQs & \draft & Verify history and prevalence\\ \hline
2. RM (I01--I03) & Discounted LP, exact MDP, fluid LP & \draft & Audit signs and example\\ \hline
3. Method (I04--I06, I09) & Affine ALP, duality, decision rule & \draft & Complete formal audit\\ \hline
4. Proof (I07, I10) & Bound hierarchy and theorem & \draft & File independent proof audit\\ \hline
5. Scale (I11) & Pricing and column generation & \draft & Verify complexity and separation\\ \hline
6. Design (I08, I12) & Benchmarks, metrics, policies & \draft & Approve factors and baselines\\ \hline
6. Results (I08, I12) & Code, tables, chart, sensitivity & \notstarted & Pass preflight and tests\\ \hline
7. Critique (I13) & Assumptions and alternatives & \draft & Link claims to evidence\\ \hline
8. Literature (I14) & Pricing lineage and FTC context & \notstarted & Verify works and primary source\\ \hline
8. Extension (I14) & Synthetic pricing application & \blocked & Await course and human gate\\ \hline
9. Conclusion & Three supported takeaways & \draft & Update after results\\
\bottomrule
\end{tabularx}
\vspace{0.7em}\par
\textbf{Key:}\quad Red \blocked\qquad Yellow \draft\qquad Amber \notstarted\qquad Green \complete
\src{User's final outline and project phase plan. Status describes presentation progress; formal requirements I01--I14 remain unverified in the project register.}
\end{frame}

% Main slide 3
\begin{frame}[t]{Roadmap from IPR 1 to the final presentation}
\footnotesize
\begin{tabularx}{\textwidth}{@{}p{0.22\textwidth}p{0.25\textwidth}p{0.25\textwidth}X@{}}
\toprule
IPR 1: scope & IPR 2: theory & IPR 3: evidence & Final: synthesis\\
\midrule
\(\bullet\) Six approved questions\newline
\(\bullet\) LP themes\newline
\(\bullet\) Verification plan\newline
\(\bullet\) Application guardrails
&
\(\bullet\) Exact and discounted LPs\newline
\(\bullet\) Fluid and affine models\newline
\(\bullet\) Bound proof\newline
\(\bullet\) Pricing formulation\newline
\(\bullet\) Experiment approval
&
\(\bullet\) Independent oracle\newline
\(\bullet\) Certified bounds\newline
\(\bullet\) Static/dynamic policies\newline
\(\bullet\) Scaling and truncation\newline
\(\bullet\) Reviewed results
&
\(\bullet\) Motivation\newline
\(\bullet\) Adelman result\newline
\(\bullet\) Measured comparisons\newline
\(\bullet\) Critique\newline
\(\bullet\) Gated application\newline
\(\bullet\) Limits
\\
\bottomrule
\end{tabularx}
\vspace{0.9em}
\textbf{Current gate:} Theory drafted. Formal audit, Windows solver preflight, code, and numerical results pending.
\medskip\par
\textbf{Final brief rule:} every numerical chart must trace to a reviewed result. The optional application cannot displace instructor-required theory or computation.
\src{Roadmap structure follows IPR 1 brief, p. 6; milestones follow the project phase plan.}
\end{frame}

% Main slide 4
\begin{frame}[t]{Research scope and current progress}
\small
\textbf{Question groups}
\begin{tabularx}{\textwidth}{@{}p{0.20\textwidth}X@{}}
\toprule
Theory (RQ1--3) & What does the affine value approximation retain? Why is its bound stronger? How do dual prices and column generation work?\\
Computation (RQ4--5) & How do size and truncation affect cost and quality? When do dynamic prices improve expected reward?\\
Application (RQ6) & Does the resource-allocation model transfer to surveillance missions?\\
\bottomrule
\end{tabularx}
\vspace{0.8em}
\begin{tabularx}{\textwidth}{@{}p{0.30\textwidth}X@{}}
\textbf{Formulations and proofs} & Source analysis drafted; formal audit artifact pending.\\
\textbf{Implementations and tests} & Exact DP, LPs, policies, and assertions pending.\\
\textbf{Experiments and replication} & Design proposed; no project run or published-table reproduction.\\
\end{tabularx}
\src{Six approved research questions and the project phase plan.}
\end{frame}

% Main slide 2 of 10
\begin{frame}[t]{Finite-horizon network revenue management}
\small
\begin{tabularx}{\textwidth}{@{}p{0.21\textwidth}X@{}}
\toprule
State & Period $t\in\{1,\ldots,T\}$ and remaining resource vector $r\in\R=\prod_{i=1}^m\{0,\ldots,c_i\}$.\\
Arrival & Class $j$ arrives with probability $p_{tj}$; no arrival has probability $1-\sum_jp_{tj}$. At most one request arrives per period.\\
Action & $u_j\in\{0,1\}$ says whether class $j$ would be accepted; $a_{ij}u_j\le r_i$ for every resource $i$ and class $j$.\\
Reward and next state & An accepted class $j$ earns $f_j$ and moves to $r-A_j$; otherwise capacity stays at $r$.\\
\bottomrule
\end{tabularx}
\vspace{0.65em}
\begin{columns}[T,onlytextwidth]
\column{0.48\textwidth}
\textbf{One-period MDP sketch}\par
\scriptsize
\[
\begin{array}{c}
\text{state }(t,r)\\[-0.1em]
\downarrow\ \text{class }j\text{ arrives with }p_{tj}\\[0.2em]
\begin{cases}
\text{accept: }(f_j,t+1,r-A_j),\\
\text{reject: }(0,t+1,r).
\end{cases}
\end{array}
\]
\column{0.48\textwidth}
\textbf{Adelman's acceptance rule}\par
\scriptsize
\[
u^*_{tj}(r)=1\quad\Longleftrightarrow\quad
r\ge A_j,\quad
f_j-\bigl[v_{t+1}(r)-v_{t+1}(r-A_j)\bigr]\ge0.
\]
Revenue today must cover the next-period value of capacity used.
\end{columns}
\notation{$t=1,\ldots,T$ period; $i=1,\ldots,m$ resource; $j=1,\ldots,n$ class; $c_i$ initial and $r_i$ remaining capacity; $a_{ij}$ units used; $A_j=(a_{ij})_i$; $p_{tj}$ arrival probability; $u_j$ accept indicator; $f_j$ fare; $v_t(r)$ optimal future reward.}
\src{Adelman (2007), pp. 648--649, \S2.1 and eqs. (1)--(2).}
\end{frame}

% Main slide 3 of 10
\begin{frame}[t]{Exact MDP and its linear program}
\small
\begin{columns}[T,onlytextwidth]
\column{0.53\textwidth}
\textbf{Bellman recursion}\par
$v_{T+1}(r)=0$ and
\[
v_t(r)=v_{t+1}(r)+\sum_{j:r\ge A_j}p_{tj}
\bigl[f_j-\Delta_{t+1,j}(r)\bigr]_+ ,
\]
\[
\Delta_{t+1,j}(r)=v_{t+1}(r)-v_{t+1}(r-A_j).
\]
For a feasible arrival, accept if
\[f_j\ge v_{t+1}(r)-v_{t+1}(r-A_j).\]
\column{0.43\textwidth}
\textbf{Bellman LP (D0)}\par
Minimize $v_1(c)$ over unrestricted $v_t(r)$ subject to
\[
v_t(r)\ge q_t(r,u)+\E[v_{t+1}(R_{t+1})\mid r,u]
\]
for every $t,r,u\in\U(r)$.
\medskip\par
Its exact dual conserves probability at every time-state pair. Complementary slackness identifies optimal actions on occupied states.
\end{columns}
\notation{$t$ period; $T$ horizon; $r$ remaining capacity; $c$ initial capacity; $A_j$ class-$j$ resource use; $p_{tj}$ arrival probability; $f_j$ fare; $v_t(r)$ optimal value; $\Delta$ opportunity cost; $[x]_+=\max(0,x)$; $u$ feasible action in $\U(r)$; $q_t$ expected reward; $R_{t+1}$ next state; $\E$ expectation.}
\src{Adelman (2007), p. 649, eqs. (1)--(3); occupancy dual independently derived in backup slide A2.}
\end{frame}

% Main slide 4 of 10
\begin{frame}[t]{Fluid LP and affine approximate LP}
\small
\begin{columns}[T,onlytextwidth]
\column{0.44\textwidth}
\textbf{Fluid LP: expected acceptances $Y_j$}
\[
\max_{Y\ge0}\ \sum_j f_jY_j\quad
\text{s.t. } AY\le c,\quad
Y_j\le\sum_t p_{tj}.
\]
Resource dual prices $\lambda_i\ge0$ give the static test
\[f_j\ge A_j^\top\lambda.\]
This solve aggregates demand over time.
\column{0.52\textwidth}
\textbf{Affine value approximation}
\[\hat v_t(r)=\theta_t+b_t^\top r.\]
Substitution in (D0) gives (D1):
\[
\begin{aligned}
\theta_t-\theta_{t+1}+(b_t-b_{t+1})^\top r
\ge \sum_j p_{tj}(f_j-b_{t+1}^\top A_j)u_j.
\end{aligned}
\]
The policy accepts a feasible class $j$ if $f_j\ge b_{t+1}^\top A_j$.
\end{columns}
\medskip\par
Within one affine solve, bid prices vary by period. Re-solving can update the trajectory after capacity changes.
\notation{$Y_j$ expected acceptances; $A=(a_{ij})$ resource-use matrix; $c$ initial capacity; $f_j$ fare; $p_{tj}$ arrival probability; $\lambda_i$ fluid resource dual; $\hat v_t$ affine value; $\theta_t$ intercept; $b_t$ resource-price vector; $r$ remaining capacity; $u_j$ accept indicator; $t$ period.}
\src{Adelman (2007), p. 650, eqs. (4)--(6), and p. 652, eqs. (14)--(17).}
\end{frame}

% Main slide 5 of 10
\begin{frame}[t]{Two upper-bound arguments}
\small
\begin{center}
{\Large $z^{LP}\ \ge\ z^{P1}\ \ge\ v_1(c)$}
\end{center}
\begin{enumerate}
\item \textbf{Affine LP versus exact MDP.} Every feasible (D1) function satisfies all Bellman inequalities. Backward induction makes it an upper bound on $v_t(r)$. Restricting the value-function family cannot lower the minimum below $v_1(c)$.
\item \textbf{Fluid LP versus affine LP.} Given feasible (P1) weights $X$, set $Y_j=\sum_{t,r,u}p_{tj}u_jX_{t,r,u}$. P1's mass and resource-moment equations imply $Y_j\le\sum_t p_{tj}$ and $AY\le c$. The two objectives agree.
\end{enumerate}
\medskip\par
\textbf{Structural result:} under $c_i>0$ and $p_{tj}>0$, a monotone optimal trajectory $b_{t,i}\ge b_{t+1,i}$ exists. An arbitrary optimal solver output need not be monotone.
\decision{Is this theorem and proof depth sufficient for the final brief?}
\notation{$z^{LP}$ fluid upper bound; $z^{P1}$ affine-LP upper bound; $v_1(c)$ exact initial MDP value; $X_{t,r,u}$ P1 state-action weight; $Y_j$ expected class-$j$ acceptance; $p_{tj}$ arrival probability; $u_j$ accept indicator; $A$ resource-use matrix; $b_{t,i}$ resource-$i$ slope; $c_i$ initial capacity.}
\src{Adelman (2007), pp. 652--655, Theorems 1--2 and eqs. (22)--(29).}
\end{frame}

% Main slide 6 of 10
\begin{frame}[t]{Published example and planned verification}
\small
\textbf{Example 1:} one resource, one seat, one class, $f=a=1$, $T=2$, and $p_1=p_2=0.5$.
\[v_2(1)=0.5,\qquad v_1(1)=0.5+0.5(0.5)=0.75.\]
\begin{center}
\begin{tabular}{@{}lcl@{}}
\toprule
Formulation & Value & Evidence status\\
\midrule
Exact MDP & 0.75 & Hand calculation; published example\\
Affine LP (P1) & 0.75 & Published example; reproduction planned\\
Fluid LP & 1.00 & Hand calculation; published example\\
\bottomrule
\end{tabular}
\end{center}
\textbf{Verification gate:} independently compare exact DP, exact LP primal/dual, fluid LP, and full affine LP on this case and a two-resource network. Test zero demand, zero capacity, terminal decisions, and infeasible arrivals.
\notation{$T$ horizon; $f$ fare; $a$ resource use; $p_t$ arrival probability; $v_t(r)$ exact value at period $t$ with remaining capacity $r$; P1 affine approximation; fluid LP deterministic relaxation.}
\src{Adelman (2007), p. 653, Example 1. Project-generated solver results are pending.}
\end{frame}

% Main slide 7 of 10
\begin{frame}[t]{Column generation and its certificate}
\small
\textbf{Primary:} start (P1) with feasible no-acceptance columns. Solve the restricted primary and obtain $\theta_t,b_t$.
\medskip\par
\textbf{Pricing:} for each period maximize reduced profit over feasible $(r,u)$:
\[
\rho_t(r,u)=\sum_j p_{tj}(f_j-b_{t+1}^{\top}A_j)u_j
-(b_t-b_{t+1})^\top r-\theta_t+\theta_{t+1}.
\]
Add a column when its reduced profit is positive, then re-solve. Adelman formulates the period problem as an integer program.
\medskip\par
\textbf{Stopping certificate (Proposition 3):}
\[
Z_{\Lambda}\le z^{P1}\le Z_{\Lambda}+\sum_{t=1}^{T}\rho_t^*,
\quad \rho_t^*=\max_{r,u}\rho_t(r,u).
\]
$Z_{\Lambda}$ alone does not certify an upper bound on the exact MDP.
\decision{How much of the pricing and column-generation algorithm must IPR 3 implement?}
\notation{$t,T$ period/horizon; $(r,u)$ state/action; $p_{tj},f_j,A_j$ arrival probability, fare, resource use; $b_t,\theta_t$ slopes/intercept; $\rho_t,\rho_t^*$ reduced profit/period maximum; $Z_\Lambda,z^{P1}$ restricted/full primary objectives.}
\src{Adelman (2007), pp. 650--651, Proposition 3; pp. 655--657, \S4.4 and eqs. (33)--(36).}
\end{frame}

% Main slide 8 of 10
\begin{frame}[t]{Proposed experiment design for IPR 3}
\small
\begin{tabularx}{\textwidth}{@{}p{0.21\textwidth}X@{}}
\toprule
Factors & Resources, horizon, load, fare spread, arrival timing, and single- versus multi-resource products. Start with exactly enumerable cases.\\
Comparators & Exact DP where tractable; fluid and certified affine bounds; fixed static and dynamic policies; separately labeled re-solving variants; proposed first-come feasible baseline.\\
Outcomes & Bound gaps, exact or estimated policy reward, paired dynamic-minus-static difference, runtime, memory, iterations, and generated columns.\\
Inference & Common arrival paths and predeclared seeds; report 95\% intervals for each policy mean and for paired differences. Use absolute gaps when exact value is zero.\\
\bottomrule
\end{tabularx}
\medskip\par
\textbf{Truncation study:} record which state, column, or horizon restriction is used and whether a valid pricing certificate survives.
\decision{Which published result and comparison protocol should IPR 3 target?}
\src{Proposed design; no project experiment has run. Adelman (2007), pp. 657--660, supplies comparison context.}
\end{frame}

% Critical appraisal before the optional application
\begin{frame}[t]{Critical appraisal for the final presentation}
\small
\begin{tabularx}{\textwidth}{@{}p{0.28\textwidth}X@{}}
\toprule
Model assumption & Question to test or disclose\\
\midrule
Known exogenous arrivals & How sensitive are policies to misspecified probabilities?\\
Nonreplenishing capacity & Does the intended application have reusable assets?\\
Affine value function & How much state dependence is lost within one solve?\\
Algorithm and scaling & Which claims are proved, certified, or only observed on test instances?\\
Published policy comparison & How much gain comes from time prices, lookahead, or re-solving?\\
\bottomrule
\end{tabularx}
\medskip\par
\textbf{Final evidence rule:} compare bounds with exact values where possible; compare policies using exact evaluation or paired simulation. Label published results separately from project measurements.
\src{Reviewer-style questions derived from Adelman (2007), \S\S2--5; empirical answers remain pending.}
\end{frame}

% Optional application
\begin{frame}[t]{Adelman and individualized pricing: a research lineage}
\small
\textbf{Distinct mechanisms}\par
Adelman (2007) estimates an \emph{internal capacity opportunity cost}. Personalized pricing also models how an individual responds to an offered price.
\medskip
\begin{tabularx}{\textwidth}{@{}p{0.25\textwidth}X@{}}
\toprule
Research thread & Candidate references supplied for verification\\
\midrule
Customer choice and related methods & Zhang and Adelman (2009); Barz et al. (2023); Hausenblas et al. (2025)\\
Individualized pricing & Jahandideh et al. (2020); Lei et al. (2023/2025)\\
Parallel personalization line & Acquisti and Varian (2005); Aydin and Ziya (2009); Ban and Keskin (2021)\\
Policy context & FTC inquiry (January 2025): consumer data in individualized pricing\\
\bottomrule
\end{tabularx}
\medskip\par
\textbf{Evidence boundary:} these are reading leads, not yet source-checked for this project. A citation trail does not establish that Adelman created surveillance-pricing systems.
\decision{Is a conceptual, source-checked connection sufficient for the final brief?}
\src{Adelman (2007), pp. 647--648; subsequent works and FTC item supplied by the researcher, pending source verification.}
\end{frame}

% Main slide 10 of 10
\begin{frame}[t]{IPR 3 gates and decisions requested}
\small
\begin{tabularx}{\textwidth}{@{}p{0.23\textwidth}X@{}}
\toprule
Next evidence gate & Source-linked proof notes; exact DP and LP agreement; fluid/static tests; full small affine LP; pricing versus exhaustive search; paired policy comparisons.\\
\textcolor{teal}{\(\bigstar\) Proof (8)} & Is the theorem and LP treatment sufficient for the final brief?\\
\textcolor{teal}{\(\bigstar\) Algorithm (10)} & How much pricing and column generation must IPR 3 implement?\\
\textcolor{teal}{\(\bigstar\) Evidence (11)} & Which published result and comparison protocol should IPR 3 target?\\
\textcolor{teal}{\(\bigstar\) Application (13)} & Is a conceptual, source-checked connection sufficient?\\
\bottomrule
\end{tabularx}
\src{ISE 631 course prompt, p. 2; project phase plan.}
\end{frame}

\appendix

\begin{frame}[t]{A1. Full transition and Bellman model}
\small
Let $\U(r)=\{u\in\{0,1\}^n:a_{ij}u_j\le r_i\ \forall i,j\}$,
$p_{t0}=1-\sum_jp_{tj}$, and $q_t(r,u)=\sum_jp_{tj}f_ju_j$. Then
\[
P_t(r'\mid r,u)=p_{t0}\mathbf1\{r'=r\}
+\sum_jp_{tj}\mathbf1\{r'=r-A_ju_j\}.
\]
Overlapping next states add their probabilities. With $v_{T+1}\equiv0$,
\[
v_t(r)=\max_{u\in\U(r)}\left\{q_t(r,u)+\sum_{r'}P_t(r'\mid r,u)v_{t+1}(r')\right\}.
\]
Initial state is $c$ in period 1. No discount factor applies to this finite-horizon model.
\notation{$t,T$ period and horizon; $r,r'$ current and next resource vectors; $c$ initial vector; $i,j,n$ resource, class, class count; $a_{ij},A_j$ class resource use; $u_j,\U(r)$ accept indicator and feasible actions; $p_{tj},p_{t0}$ arrival/no-arrival probabilities; $f_j$ fare; $q_t$ expected reward; $P_t$ transition; $v_t$ optimal value; $\mathbf1$ indicator.}
\src{Adelman (2007), pp. 648--649, eq. (1); transition notation derived.}
\end{frame}

\begin{frame}[t]{A2. Exact finite-horizon Bellman LP and occupancy dual}
\footnotesize
\textbf{Primal (D0):} $\min v_1(c)$, with $v_t(r)\in\mathbb R$, $v_{T+1}=0$, and
\[v_t(r)-\sum_{r'}P_t(r'\mid r,u)v_{t+1}(r')\ge q_t(r,u)
\quad\forall t,r\in\R_t,u\in\U(r).\]
\textbf{Derived dual:} $x_t(r,u)\ge0$ and
\[
\max\sum_{t,r,u}q_t(r,u)x_t(r,u)
\]
subject to
\[
\sum_u x_1(c,u)=1,\qquad
\sum_u x_t(r,u)=\sum_{s,u}P_{t-1}(r\mid s,u)x_{t-1}(s,u)
\quad\forall t=2,\ldots,T,\ r\in\R_t.
\]
Strong duality gives the exact initial value. Positive $x_t(r,u)$ makes its Bellman inequality tight.
\notation{$t,T$ period/horizon; $r,s$ states in $\R_t$; $u\in\U(r)$ feasible action; $c$ initial state; $q_t$ expected reward; $P_t$ transition probability; $v_t$ free value variable; $x_t$ nonnegative state-action occupancy; $v_{T+1}=0$ terminal value.}
\src{Adelman (2007), p. 649, eq. (3); dual independently derived, not labeled P0 in the paper.}
\end{frame}

\begin{frame}[t]{A3. Instructor's discounted MDP LP}
\small
For finite states $S$, actions $A(s)$, discount $0\le\gamma<1$, and initial weight $\nu\ge0$:
\[
\begin{aligned}
\min_{V\in\mathbb R^{S}}\quad&\sum_s\nu(s)V(s)\\
\text{s.t.}\quad&V(s)-\gamma\sum_{s'}P(s'\mid s,a)V(s')\ge r(s,a)
&&\forall s,a\in A(s).
\end{aligned}
\]
The occupancy dual maximizes $\sum_{s,a}r(s,a)x(s,a)$ over $x\ge0$ with
\[
\sum_a x(s,a)-\gamma\sum_{s',a'}P(s\mid s',a')x(s',a')=\nu(s)
\quad\forall s.
\]
$x$ counts $\nu$-weighted discounted visits. Complementary slackness identifies optimal actions on occupied states. With $\nu(s)>0$ for all $s$, the minimizing value vector is $V^*$ everywhere.
\notation{$S$ state set; $A(s)$ feasible actions; $s,s'$ states; $a$ action; $\gamma$ discount; $\nu$ initial state weights; $V,V^*$ candidate/optimal values; $P$ transition probability; $r(s,a)$ reward; $x(s,a)$ discounted occupancy.}
\src{ISE 631 instructor prompt, p. 1; dual and full-support qualification derived.}
\end{frame}

\begin{frame}[t]{A4. Fluid primal, dual, and policy distinction}
\small
Let $d_j=\sum_t p_{tj}$. The fluid primal is
\[
\max_{Y\ge0}\ \sum_jf_jY_j\quad
\text{s.t. } AY\le c,\quad Y_j\le d_j\ \forall j.
\]
Its dual is
\[
\min_{\lambda,\mu\ge0}\ c^\top\lambda+d^\top\mu
\quad\text{s.t. } A_j^\top\lambda+\mu_j\ge f_j\ \forall j.
\]
Capacity shadow price $\lambda_i$ is nonnegative. Complementary slackness characterizes primal-dual optima; an online rule accepting when $f_j\ge A_j^\top\lambda$ is a policy constructed from these prices, with separate performance to be evaluated.
\notation{$t$ period; $j$ class; $d_j=\sum_t p_{tj}$ expected demand; $Y_j$ expected acceptances; $f_j$ fare; $A=(a_{ij})$ resource use; $c$ initial capacity; $\lambda_i$ resource dual; $\mu_j$ demand-cap dual.}
\src{Adelman (2007), p. 650, eqs. (4)--(6) and displayed dual.}
\end{frame}

\begin{frame}[t]{A5. Full affine Bellman LP (D1)}
\small
With $\R_1=\{c\}$, $\R_t=\R$ for $t=2,\ldots,T$, and
$\theta_{T+1}=0$, $b_{T+1}=0$, the raw affine coefficients
$\theta_t,b_{t,i}\in\mathbb R$ are free. The program is
\[
\min_{\theta,b}\quad \theta_1+b_1^\top c
\]
subject to
\[
\begin{gathered}
\theta_t-\theta_{t+1}+(b_t-b_{t+1})^\top r
\ge \sum_jp_{tj}(f_j-b_{t+1}^\top A_j)u_j\\
\forall t=1,\ldots,T,\quad r\in\R_t,\quad u\in\U(r).
\end{gathered}
\]
Theorem 2 permits a nonnegative monotone optimal selection only under its stated positive-capacity and positive-arrival assumptions.
\notation{$t=1,\ldots,T$ period; $r\in\R_t$ remaining state; $c$ initial capacity; $u\in\U(r)$ feasible acceptance vector; $\theta_t$ free intercept; $b_t$ free resource slopes; $p_{tj}$ arrival probability; $f_j$ fare; $A_j$ class-$j$ resource use.}
\src{Adelman (2007), p. 652, eqs. (14), (16)--(17); Theorem 2, pp. 654--655.}
\end{frame}

\begin{frame}[t]{A6. Affine dual and the bound proof}
\footnotesize
For (P1), $X_t(r,u)\ge0$ maximizes $\sum_{t,r,u}q_t(r,u)X_t(r,u)$ subject to
\[
\sum_{r,u}X_t(r,u)=1\quad\forall t,
\]
\[
\sum_{r,u}r_iX_1(r,u)=c_i,
\quad
\sum_{r,u}r_iX_t(r,u)=\sum_{s,u}\left(s_i-\sum_jp_{t-1,j}a_{ij}u_j\right)X_{t-1}(s,u)
\quad\forall i,t>1.
\]
These are moments, not statewise flows. Set $Y_j=\sum_{t,r,u}p_{tj}u_jX_t(r,u)$.
Mass gives $Y_j\le d_j$. Telescoping resource moments gives
\[
c_i=\sum_{t<T,r,u,j}p_{tj}a_{ij}u_jX_t(r,u)
+\sum_{r,u}r_iX_T(r,u)\ge\sum_ja_{ij}Y_j,
\]
since $r_i\ge\sum_jp_{Tj}a_{ij}u_j$. Therefore every P1 objective is fluid feasible. D1 feasibility separately gives $z^{P1}\ge v_1(c)$.
\notation{$t,T$ period/horizon; $i,j$ resource/class; $r,s$ remaining states; $u_j$ accept indicator; $X_t(r,u)$ P1 weight; $q_t$ expected reward; $p_{tj}$ arrival probability; $a_{ij}$ resource use; $c_i$ initial capacity; $Y_j,d_j$ expected acceptances/arrivals; $z^{P1},v_1(c)$ affine bound/exact value.}
\src{Adelman (2007), pp. 652--653, eqs. (18)--(24), Theorem 1.}
\end{frame}

\begin{frame}[t]{A7. Monotone prices: scope of Theorem 2}
\small
Under $c_i>0$ and $p_{tj}>0$ for every $i,t,j$, an optimal D1 solution exists with
\[b_{t,i}\ge b_{t+1,i}\ge0,\qquad \theta_t\ge\theta_{t+1}\ge0.\]
The proof uses the first period in which a resource appears in positive P1 support, complementary slackness, and changes to the dual prices that preserve feasibility and optimality.
\medskip\par
\textbf{Why ``exists'' matters:} Example 1 also admits $\theta_2=0$, $b_2=0.5$, $\theta_1=0.75$, $b_1=0$. Its D1 objective is $0.75$, but $b_1<b_2$. Zero-demand or zero-capacity tests lie outside the theorem's stated assumptions.
\notation{$t$ period; $i$ resource; $j$ class; $c_i$ initial capacity; $p_{tj}$ arrival probability; $b_{t,i}$ affine resource price; $\theta_t$ affine intercept; D1/P1 are the affine primal-dual pair.}
\src{Adelman (2007), pp. 653--655, eqs. (25)--(29), Theorem 2 and Corollary 1; counterexample derived.}
\end{frame}

\begin{frame}[t]{A8. Pricing, reduced costs, and stopping}
\footnotesize
For fixed $t>1$, define $w_j=p_{tj}(f_j-b_{t+1}^\top A_j)$ and $d_i=b_{t,i}-b_{t+1,i}$. Adelman's pricing problem is
\[
\max_{r,u}\ \sum_jw_ju_j-\sum_id_ir_i-\theta_t+\theta_{t+1}
\quad\text{s.t. }a_{ij}u_j\le r_i,\quad u_j\in\{0,1\},\quad 0\le r_i\le c_i.
\]
The paper uses an integer program and generates a batch of improving columns. With global period maxima $\rho_t^*$, Proposition 3 yields $z^{P1}\le Z_\Lambda+\sum_t\rho_t^*$; if $Z_\Lambda>0$, $\sum_t\rho_t^*/Z_\Lambda\le\varepsilon$ is a sufficient relative stopping rule.
\medskip\par
\textbf{Complexity caveat:} A derived minimum-cut reduction for this exact pricing structure is being proof-checked. The paper itself does not claim NP-hardness or polynomial complexity.
\notation{$t,T$ period/horizon; $i,j$ resource/class; $w_j$ net expected class profit; $d_i$ price change; $r_i,c_i$ remaining/initial capacity; $a_{ij}$ resource use; $u_j$ accept indicator; $\theta_t,b_t$ intercept/slopes; $\rho_t^*$ max reduced profit; $Z_\Lambda$ restricted-primary value; $z^{P1}$ full value; $\varepsilon$ tolerance.}
\src{Adelman (2007), pp. 650--651, Proposition 3, and pp. 655--657, eqs. (33)--(36), Figure 2.}
\end{frame}

\begin{frame}[t]{A9. Proposed experiment matrix}
\footnotesize
\begin{tabularx}{\textwidth}{@{}p{0.18\textwidth}p{0.25\textwidth}X@{}}
\toprule
Block & Variation & Evidence to record\\
\midrule
Verification & Example 1; 2-resource case; edge cases & DP, exact LP and dual, fluid, full affine LP; values and residuals\\
Bound strength & Size, load, timing, fare & Exact $V^*$ when feasible; fluid and certified affine bounds\\
Policy value & Fixed and re-solved policies & Tiny-case exact value; policy and paired-difference 95\% intervals\\
Computation & Enumeration vs.\ columns & Runtime, memory, iterations, generated columns, pricing gaps\\
Truncation & Columns, states, horizon & Time saved, omitted constraints, surviving certificates\\
\bottomrule
\end{tabularx}
\medskip\par
Factor levels, pilot size, final replication precision, seeds, and published replication target require approval before Phase 4 runs.
\src{Project experiment design proposal; no project experimental results.}
\end{frame}

\begin{frame}[t]{A10. Instructor requirements I01--I07}
\footnotesize
\begin{tabularx}{\textwidth}{@{}p{0.08\textwidth}p{0.15\textwidth}p{0.16\textwidth}X@{}}
\toprule
ID & Final part & Register & Evidence needed next\\
\midrule
I01 & 2--4 & Planned & Discounted Bellman LP, occupancy dual, complementary slackness audit\\
I02 & 2 & Planned & Exact network MDP and state-growth analysis\\
I03 & 2--3 & Planned & Fluid primal/dual and static-policy validation\\
I04 & 3--5 & Planned & Source-matched affine ALP and algorithm audit\\
I05 & 3 & Planned & Paper method and solution approach, source-linked\\
I06 & 3--5 & Planned & Correct LP course-concept mapping\\
I07 & 4 & Planned & Principal theorem and independently reviewed proof\\
\bottomrule
\end{tabularx}
\medskip\par
Draft formulas in this brief are presentation material. The register remains planned until evidence paths and a separate reviewer are recorded.
\src{\texttt{requirements.csv}, I01--I07; instructor prompt, pp. 1--2.}
\end{frame}

\begin{frame}[t]{A11. Instructor requirements I08--I14}
\footnotesize
\begin{tabularx}{\textwidth}{@{}p{0.08\textwidth}p{0.15\textwidth}p{0.16\textwidth}X@{}}
\toprule
ID & Final part & Register & Evidence needed next\\
\midrule
I08 & 6 & Planned & Reproducible fluid, static, and dynamic policy tests\\
I09 & 3--4 & Planned & Approximation benefit with verified example\\
I10 & 4 & Planned & Full bound-dominance derivation and audit\\
I11 & 5 & Planned & Pricing formulation and qualified complexity account\\
I12 & 6 & Planned & Policy, exact MDP, and fluid gap results\\
I13 & 7 & Planned & Evidence-based critique and alternatives\\
I14 & 1, 5, 8 & Planned & Approved independent RQs linked to final analysis\\
\bottomrule
\end{tabularx}
\medskip\par
I14 has approved questions on file, but its requirement register is not reviewed. I08 and I12 have no project numerical results yet.
\src{\texttt{requirements.csv}, I08--I14; instructor prompt, p. 2.}
\end{frame}

\end{document}
